Aggregate queries compress many records into a single number, but the
number of contributing records can obscure dependence on a few source
tuples. In a join, deleting one connector tuple can remove many aggregate
inputs at once. A count supported by thousands of derivations may therefore
cross a decision threshold after only a handful of source deletions.
This motivates a source-level question: how many permitted deletions are
needed to change the quantitative conclusion?

Small-deletion robustness is also studied in statistical analysis, where
removing a few observations can change an estimate or its
interpretation~\cite{broderick2020dropping,huang2025dropping}. Database
joins introduce an additional dependency: one source tuple can participate
in many output derivations. Local sensitivity measures the effect of a
single update~\cite{tao2020sensitivity}, while ranking stability studies
how changes to an item's values affect its
position~\cite{campbellLocalStability}. We instead fix an aggregate-answer
threshold and minimize the number of source deletions needed to cross it.

We call an optimal deletion set a \emph{minimum instability witness}
(\MIW). The threshold is directional: downward for \agg{COUNT},
nonnegative \agg{SUM}, and \agg{MAX}; upward for \agg{MIN}; and either
chosen direction for \agg{AVG}. Designated deletability lets an analyst
vary uncertain records while treating reference data as immutable.
Minimum cardinality is essential: an inclusion-minimal set may contain
many more deletions than necessary.

\paragraph{Relationship to deletion propagation.}
For full self-join-free conjunctive queries over duplicate-free relations,
the COUNT objective coincides with aggregated deletion propagation (ADP).
Hu et al.~\cite{hu2021adp} give logarithmic and relation-count
approximations in that shared setting. Classical view deletion and its
source/view side-effect objectives provide the broader
context~\cite{buneman2002deletion,kimelfeld2011deletion}; generalized
deletion propagation also supplies ILP methods for combined deletion and
preservation requirements~\cite{makhija2025ilp}. We build on this foundation
to explain how numerical thresholds and nonempty aggregate inputs determine
the optimization structure across five operators.

\paragraph{Three aggregate regimes.}
COUNT and nonnegative SUM yield weighted coverage: deleting a tuple removes
the nonnegative contributions of its previously surviving derivations.
MAX and MIN instead require hitting every derivation on the wrong side of
the target while preserving an acceptable survivor. Their immediate
numerical change can conceal progress until all competing extrema vanish.
AVG changes both numerator and denominator, giving a signed threshold
constraint. These regimes determine which exact and approximate methods
apply, and what the provenance representation must support.

\paragraph{Contributions.}
\begin{itemize}
\item We formalize directional \MIW\ under bag semantics with immutable
tuples and explicit empty-input conventions. We establish fixed-query
hardness for all five aggregates and exact algorithms for singleton
supports, showing that reusing a single deletable relation can move the
query class from uniform tractability to NP-complete cases
(Sections~\ref{sec:problem}--\ref{sec:complexity}).
\item We give coverage bounds for COUNT and SUM with explicit precision,
a complete extrema feasibility test, and a $\min\{r,1+\ln m\}$ extrema approximation,
where $r$ bounds support size and $m$ counts critical derivations. For AVG,
we identify a sign-separated coverage subclass alongside the general
signed formulation (Sections~\ref{sec:theory}--\ref{sec:algorithms}).
\item We specify a unified exact ILP and exact marginal computation through
incremental supports, lazy query oracles, and multilinear circuits. We
connect the derivative condition to source reuse and distinguish current
deletion gains from coalition-based attribution
(Section~\ref{sec:algorithms}).
\item We implement and compare these execution choices in a reproducible
TPC-H-derived pilot for COUNT, SUM, and AVG, validating every returned
witness by residual SQL and exact arithmetic
(Section~\ref{sec:experiments}).
\end{itemize}
Section~\ref{sec:related} situates these results in provenance, attribution,
robustness, and repair research. Section~\ref{sec:conclusion} discusses
extensions beyond the current model and evaluation.
